\begin{problem}[第35届全国中学生物理竞赛决赛][半球与小球动力学]{96}
    如图，半径为 $R$ 、质量为 $M$ 的半球静置于光滑水平桌面上，在半球顶点上有一质量为 $m$ 、半径为 $r$ 的匀质小球。某时刻，小球受到微小的扰动后由静止开始沿半球表面运动。在运动过程中，小球相对于半球的位置由角位置 $\theta$ 描述， $\theta$ 为两球心的连线与竖直方向之间的夹角。已知小球绕其对称轴的转动惯量为 $\frac{2}{5} mr^2$ ，小球与半球之间的动摩擦因数为 $\mu$ ，假定最大静摩擦力等于滑动摩擦力。重力加速度大小为 $g$ 。
    \begin{subproblem}
        小球开始运动后在一段时间内做纯滚动，求在此过程中，当小球的角位置为 $\theta_{1}$ 时，半球运动的速度大小 $V_{M}(\theta_{1})$ 和加速度大小 $a_{M}(\theta_{1})$ ；
    \end{subproblem}
    \begin{subproblem}
        当小球纯滚动到角位置 $\theta_{2}$ 时开始相对于半球滑动，求 $\theta_{2}$ 所满足的方程（可用半球速度大小 $V_{M}\left(\theta_{2}\right)$ 和加速度大小 $a_{M}\left(\theta_{2}\right)$ 以及题给条件表示）；
    \end{subproblem}
    \begin{subproblem}
        当小球刚好运动到角位置 $\theta_{3}$ 时脱离半球，求此时小球质心相对于半球运动速度的大小 $v_{m}(\theta_{3})$ 。
    \end{subproblem}
\end{problem}